\section{Conclusion}

Au cours de cette étude, nous avons mis en place et optimisé un dispositif optique permettant d'analyser la dynamique d'une figure de speckle thermique. En éclairant une suspension de nanoparticules dans un milieu très visqueux avec un laser à 532 nm, et en capturant la lumière diffusée à l'aide d'une caméra monochrome, nous avons pu suivre l'évolution du mouvement brownien en fonction de la température. Le calcul des fonctions d'autocorrélation temporelle de l'intensité nous a permis d'extraire le temps de relaxation et d'en déduire le coefficient de diffusion.

Les tests initiaux à température ambiante ont confirmé la fiabilité du montage en mesurant un diamètre de particule de 93 nm, une valeur en adéquation avec les données théoriques. Un résultat technique déterminant a été l'ajustement de l'orientation de la cellule : en tournant légèrement la cuve sans déplacer la caméra, nous avons éliminé les reflets parasites du polystyrène, rendant le signal de diffusion bien plus net. 

Enfin, les mesures montrent que le coefficient de diffusion augmente de façon cohérente avec la température, suivant la loi de Stokes-Einstein. L'absence de différence entre les deux protocoles de chauffe prouve qu'il n'y a pas eu de convection, la forte viscosité du mélange ayant stabilisé le milieu.

Concernant notre problématique initiale, qui consistait à vérifier si un chauffage dynamique perturbait la mesure par l'apparition de mouvements de convection, la comparaison de nos deux protocoles a apporté une réponse claire. Nous n'avons observé aucune décorrélation anormale ni accélération soudaine du speckle lors du chauffage continu par rapport au chauffage par paliers. L'absence de ces flux hydrodynamiques macroscopiques s'explique par la viscosité extrêmement élevée de notre mélange à 90 % de glycérol.

Nous avons donc démontré que dans un milieu aussi visqueux, le mouvement brownien reste le seul moteur de la dynamique, indépendamment du mode de chauffe choisi.